\section{Tổ hợp}
\begingroup\small
\begin{multicols}{2}
\subsection{Tổ hợp và đếm cách phân phối}
Với $0\le k\le n$, số cách chọn $k$ trong $n$ phần tử phân biệt,
không xét thứ tự, là
\[
\binom nk=\frac{n!}{k!(n-k)!},\qquad
\binom nk=\binom{n-1}{k-1}+\binom{n-1}k.
\]
Quy ước $\binom nk=0$ nếu $k<0$ hoặc $k>n$.
Số hoán vị của đa tập có $n_i$ phần tử loại $i$ và
$\sum_i n_i=n$ là $n!/\prod_i n_i!$.

Số nghiệm nguyên không âm của $x_1+\cdots+x_m=s$ là
$\binom{s+m-1}{m-1}$ khi $m\ge1,s\ge0$ (đặt $m-1$ vách ngăn).
Nếu $x_i\ge l_i$, thay $y_i=x_i-l_i$.
Nếu thêm $x_i\le u_i$, đặt $t=s-\sum_i l_i$,
$b_i=u_i-l_i\ge0$. Số nghiệm là
\[
\sum_{J\subseteq[m]}(-1)^{|J|}
\binom{t-\sum_{j\in J}(b_j+1)+m-1}{m-1}.
\]
Đây là bao hàm--loại trừ trên các điều kiện xấu $y_i>b_i$.
Nếu $t<0$, đáp án bằng $0$.

\subsection{Số Catalan}
Đặt $C_0=1$. Tách dãy ngoặc đúng theo cặp ngoặc đầu tiên:
\[
C_n=\sum_{i=0}^{n-1}C_iC_{n-1-i}
=\binom{2n}{n}-\binom{2n}{n+1}
=\frac{1}{n+1}\binom{2n}{n}.
\]
$C_n$ đếm dãy ngoặc đúng $n$ cặp, đường đi lưới từ
$(0,0)$ tới $(n,n)$ với mọi tiền tố giữ $y\le x$,
và số cách tam giác hóa đa giác lồi $n+2$ đỉnh (với $n\ge1$).
Hiệu hai hệ số nhị thức suy ra từ nguyên lý phản xạ.
Khi tính modulo $M$, không chia cho $n+1$ nếu nó không khả nghịch;
dùng hiệu hai hệ số nhị thức hoặc cách khác thích hợp.

\subsection{Bao hàm--loại trừ}
Với các tập hữu hạn $A_1,\ldots,A_k$,
\[
\left|\bigcup_{i=1}^k A_i\right|=
\sum_{\varnothing\ne J\subseteq[k]}(-1)^{|J|+1}
\left|\bigcap_{j\in J}A_j\right|.
\]
Một phần tử nằm trong đúng $r\ge1$ tập được tính với hệ số
$\sum_{j=1}^r(-1)^{j+1}\binom rj=1$.
Để đếm đối tượng không mắc điều kiện xấu nào, lấy toàn bộ
trừ số đối tượng thuộc hợp. Phương pháp trực tiếp có $2^k$
giao, chỉ dùng khi số điều kiện đủ nhỏ hoặc có thể gom các
tập con theo một đại lượng chung.

\subsection{Định lý Burnside}
Nhóm hữu hạn $G$ tác động lên tập cấu hình hữu hạn $X$; hai
cấu hình tương đương nếu một phép trong $G$ biến cấu hình này
thành cấu hình kia. Số lớp tương đương:
\[
|X/G|=\frac{1}{|G|}\sum_{g\in G}|\operatorname{Fix}(g)|,\qquad
\operatorname{Fix}(g)=\{x:g(x)=x\}.
\]
Đếm các cặp $(g,x)$ với $g(x)=x$ theo $g$ cho tổng ở tử;
theo định lý quỹ đạo--ổn định, mỗi lớp đóng góp $|G|$ cặp.
Nếu $g$ hoán vị $n$ vị trí và tô bằng $q$ màu, cấu hình cố định
phải có cùng màu trên mỗi chu trình, nên
$|\operatorname{Fix}(g)|=q^{c(g)}$. Tô vòng $n$ ô, chỉ đồng nhất
theo phép quay, cho
\[
\frac1n\sum_{r=0}^{n-1}q^{\gcd(n,r)}.
\]
Nếu còn đồng nhất qua phép lật, phải cộng các phép lật vào nhóm.
Khi tính modulo, phép chia cho $|G|$ bằng nghịch đảo chỉ hợp lệ
nếu $\gcd(|G|,M)=1$.

\paragraph{Hệ số nhị thức modulo số nguyên tố (C++17).}
Yêu cầu $p$ nguyên tố và $0\le N<p$; sau tiền xử lý $O(N)$,
mỗi truy vấn $O(1)$. Hàm Catalan cần $2n\le N$.
\cppsnippet{code/combinatorics/binomial.hpp}

\subsection{Nội suy Lagrange}
Trên trường $\mathbb F_p$ với $p$ nguyên tố, cho $k+1$ điểm
$(x_i,y_i)$ có các $x_i$ đôi một khác nhau modulo $p$.
Đa thức duy nhất $P$ bậc không quá $k$ đi qua các điểm đó là
\[
P(t)=\sum_{i=0}^{k} y_i
\prod_{\substack{0\le j\le k\\j\ne i}}
\frac{t-x_j}{x_i-x_j}\pmod p.
\]
Thật vậy, hạng ứng với $i$ bằng $y_i$ tại $t=x_i$ và bằng $0$
tại các mốc còn lại; hiệu của hai đa thức bậc không quá $k$
có $k+1$ nghiệm nên bằng $0$. Phép chia modulo chỉ hợp lệ khi
mẫu khác $0$ modulo $p$; tính nghịch đảo $a^{p-2}\bmod p$.

Trường hợp thường gặp: biết $y_i=P(i)$ với $0\le i\le k<p$.
Khi đó
\[
\prod_{j\ne i}(i-j)=(-1)^{k-i}i!(k-i)!.
\]
Tích tiền tố và hậu tố của $(t-j)$ cho mọi tử số trong $O(k)$;
nghịch đảo các giai thừa lấy từ một lần lũy thừa nhanh.
Mỗi lần tính $P(t)$ tốn $O(k+\log p)$ thời gian và $O(k)$ bộ nhớ.
Nếu $t\bmod p$ trùng một mốc, trả ngay $y_i$.
Điều kiện bậc không quá $k$ là thiết yếu: chỉ $k+1$ giá trị
không xác định được đa thức bậc cao hơn.

\paragraph{Nội suy trên các mốc $0,1,\ldots,k$ (C++17).}
\cppsnippet{code/combinatorics/lagrange.hpp}
\end{multicols}
\endgroup
